\documentclass[11pt]{article}
\usepackage[margin=1in]{geometry}
\usepackage{amsmath,amssymb}
\usepackage[hidelinks]{hyperref}

\title{The determinant does not determine mean Lee grading\\[2pt]
  \large Disproof of Conjecture 00000000617}
\author{}
\date{}

\begin{document}
\maketitle

Over $\mathbb Q$, Lee's theorem gives dimension $2^c$ for a link with
$c$ components, with a canonical generator for each orientation
of the components~\cite[Theorem~4.2 and \S4.4.3]{lee}. The proposed mean
formula fails for the quantum filtration and for the homological grading.

\section*{The quantum filtration on the canonical basis}

Let $K$ be the positive trefoil, the closure of $\sigma_1^3$ on two
strands, and let $\overline K$ be its mirror. The positive diagram has
three crossings. Its oriented smoothings replace the crossings by two
parallel strands, which close to two Seifert circles. Rasmussen's formula
for a positive diagram~\cite[\S5.2]{rasmussen} gives
\[
  s(K)=3-2+1=2,
  \qquad s(\overline K)=-s(K)=-2,
\]
where the mirror identity is Proposition~3.9 of the same paper. By
Corollary~3.6, each of the two canonical orientation generators of a knot
has filtration level $s_{\min}=s-1$. Their mean levels are therefore $1$
on $K$ and $-3$ on $\overline K$. The two associated-graded levels are
$s-1$ and $s+1$ by Proposition~3.3, so their means under that convention
are $2$ and $-2$, respectively.

The trefoil Alexander polynomial is $t^2-t+1$ up to multiplication by
$\pm t^a$~\cite[Example~B.1]{nicolaescu}. Mirroring replaces $t$ by
$t^{-1}$, up to the same units. Thus both determinants, defined as the
absolute Alexander value at $-1$, are $3$. The equal determinants and
unequal means refute a determinant-only formula in either filtration
convention.

\section*{The homological grading}

Let $H$ be the positive Hopf link, the closure of $\sigma_1^2$ with the
braid orientation, and give its mirror $\overline H$ the reflected
orientation. Each diagram has two intercomponent crossings, both positive
for $H$ and both negative for $\overline H$. In the usual normalized
grading, the canonical generator for an orientation has degree $r-n_-$,
where $r$ counts the $1$-resolutions in its oriented smoothing and $n_-$
counts negative crossings of the reference diagram. The oriented
smoothing is a $1$-resolution at a negative crossing and a $0$-resolution
at a positive crossing. This is the resolution calculation in the proof
of Lee's homological-degree formula~\cite[Proposition~4.3]{lee}.

Reversing neither component or both components leaves all crossing signs
unchanged. Reversing exactly one component changes both signs. Thus the
four canonical degrees are $0,0,2,2$ for $H$ and $0,0,-2,-2$ for
$\overline H$, with means $1$ and $-1$. A Hopf link bounds a once-twisted
annulus. In a core basis its Seifert matrix is $[1]$ or $[-1]$, with the
sign reversed by mirroring, because the linking number of a core and a
push-off is the signed full twist. The determinant
$\lvert\det(V+V^{\mathsf T})\rvert$ is consequently $2$ for both links.
These equal determinants also preclude a formula for the homological mean.

\section*{Lean verification}

The Lean 4 project computes the crossing and orientation data of these
two-strand closures, the resulting grading levels from the formulas
above, and the determinant evaluations. It checks the different means
and rules out a function of the determinant for each of the three
conventions. Run \texttt{lake build} in the \texttt{lean} directory to
verify the computations.

\begin{thebibliography}{3}
\bibitem{lee}
E.~S. Lee, \emph{An endomorphism of the Khovanov invariant},
Advances in Mathematics \textbf{197} (2005), 554--586.
\href{https://arxiv.org/abs/math/0210213}{arXiv:math/0210213}.

\bibitem{rasmussen}
J.~Rasmussen, \emph{Khovanov homology and the slice genus},
Inventiones Mathematicae \textbf{182} (2010), 419--447.
\href{https://arxiv.org/abs/math/0402131}{arXiv:math/0402131}.

\bibitem{nicolaescu}
L.~I. Nicolaescu, \emph{The Reidemeister Torsion of 3-Manifolds},
de Gruyter, 2003, \S B.1.
\href{https://www3.nd.edu/~lnicolae/Torsion.pdf}{Author's copy}.
\end{thebibliography}

\end{document}
